%%%%%%%%%%%%
%
% $Autor: Wings $
% $Datum: 2019-03-05 08:03:15Z $
% $Pfad: SafetyGuidelines.tex $
% $Version: 4250 $
% !TeX spellcheck = en_GB/de_DE
% !TeX encoding = utf8
% !TeX root = manual 
% !TeX TXS-program:bibliography = txs:///biber
%
%%%%%%%%%%%%

\chapter{Sicherheitshinweise}


\section{Allgemeines}
ELMAR ist eine CNC-gesteuerte Heißdraht-Schneidanlage zur Bearbeitung von Schaumstoffen. Der Betrieb ist ausschließlich gemäß der vorgesehenen Verwendung und unter Einhaltung der folgenden Sicherheitsvorschriften zulässig.

\section{Bestimmungsgemäße Verwendung}
ELMAR ist ausschließlich zum Schneiden von geeigneten Schaumstoffen vorgesehen. Eine Verwendung für andere Materialien oder Anwendungen gilt als nicht bestimmungsgemäß. Umbauten oder Veränderungen an ELMAR sind unzulässig.


\subsection{Elektrische Gefährdung}
Es besteht Gefahr durch elektrischen Strom bei beschädigten Leitungen oder unsachgemäßem Anschluss.

\textbf{Schutzmaßnahmen:}
\begin{itemize}
	\item ELMAR vor Arbeiten spannungsfrei schalten
	\item Elektrische Komponenten regelmäßig prüfen
	\item Feuchtigkeit vermeiden
\end{itemize}

\subsection{Thermische Gefährdung}
Der Heizdraht erreicht hohe Temperaturen. Es besteht Verbrennungsgefahr.

\textbf{Schutzmaßnahmen:}
\begin{itemize}
	\item Heizdraht während und nach dem Betrieb nicht berühren
	\item Abkühlzeit einhalten
	\item Heiße Bereiche kennzeichnen
\end{itemize}

\subsection{Mechanische Gefährdung}
Bewegliche Komponenten können Quetsch- und Scherstellen verursachen.

\textbf{Schutzmaßnahmen:}
\begin{itemize}
	\item Nicht in den Bewegungsbereich greifen
	\item Sicherheitsabstand einhalten
	\item Keine lose Kleidung tragen
\end{itemize}

\subsection{Brand- und Rauchentwicklung}
Beim Schneiden können Rauch und Dämpfe entstehen. Es besteht Brandgefahr.

\textbf{Schutzmaßnahmen:}
\begin{itemize}
	\item Nur in gut belüfteten Bereichen betreiben
	\item Abstand zu brennbaren Materialien einhalten
	\item Feuerlöscher bereithalten
\end{itemize}

\section{Betrieb}
ELMAR darf nur von eingewiesenem Personal betrieben werden. Vor jeder Inbetriebnahme ist eine Sichtprüfung durchzuführen. ELMAR darf nicht unbeaufsichtigt betrieben werden. Bei Störungen ist ELMAR sofort außer Betrieb zu nehmen.

\section{Wartung}
Wartungsarbeiten dürfen nur im spannungsfreien Zustand durchgeführt werden. Verschleißteile sind regelmäßig zu prüfen und bei Bedarf zu ersetzen. Beschädigte Bauteile dürfen nicht weiterverwendet werden.

\section{Notfallmaßnahmen}
Im Gefahrenfall ist ELMAR sofort vom Stromnetz zu trennen. Bei Brand sind geeignete Löschmittel zu verwenden. Verletzte Personen sind umgehend zu versorgen.



